% \newif\ifdraft \drafttrue
% \drafttrue
\documentclass[11pt]{article}
\usepackage{amssymb,amsmath,amsthm,amsfonts}
\usepackage{jinshuomacros}
\usepackage{algorithm}
\usepackage[noend]{algpseudocode}
\usepackage{thmtools}
\usepackage{mathrsfs}
\usepackage[usenames,dvipsnames]{color}
\usepackage{tikz}
% \usepackage{bbm}
% \usepackage[pdftex]{hyperref}
\usepackage{hyperref}

\usepackage{graphicx}
\usepackage{bibentry}
\usepackage{cleveref}
\usepackage{fullpage}

\newif\ifhideproofs
% \hideproofstrue %uncomment to hide proofs

\ifhideproofs
\usepackage{environ}
\NewEnviron{hide}{}
\let\proof\hide
\let\endproof\endhide
\fi
% \newcommand{\crefrangeconjunction}{--}
\renewcommand{\N}{\mathcal{N}}
\renewcommand{\P}{\mathbb{P}}
\newcommand{\e}{\mathrm{e}}
\newcommand{\f}{\boldsymbol{f}}
\newcommand{\DCT}{{dominated convergence theorem}}
%%\mathrm{Proc}
\newcommand{\Proc}{\mathrm{Proc}}
\newcommand{\Renyi}{R\'enyi}
\renewcommand{\algorithmicrequire}{\textbf{Input:}}
\renewcommand{\algorithmicensure}{\textbf{Output:}}

% \newcommand{\djs}[1]{\ifdraft \textcolor{blue}{[Jinshuo: #1]}\fi}
% \newcommand{\ar}[1]{\ifdraft \textcolor{red}{[Aaron: #1]}\fi}
% \newcommand{\wjs}[1]{\ifdraft \textcolor{purple}{[Weijie: #1]}\fi}
% \newcommand{\weijie}[1]{\ifdraft \textcolor{red}{[Weijie: #1]}\fi}

\newcommand{\tv}{d_{\textnormal{TV}}}
% \newcommand{\TV}{{d_{\textnormal{TV}}}}
\newcommand{\Sf}{f^{\mathrm{S}}}
%\DeclareMathOperator{\TV}{d_{TV}}
		\newcommand{\TV}{\mathrm{TV}}

\nobibliography* % See https://tex.stackexchange.com/questions/49048/how-to-cite-one-bibentry-in-full-length-in-the-body-text
% for an explanation why this line is necessary

%     If the article includes commutative diagrams, ...
%\usepackage[cmtip,all]{xy}

% \copyrightinfo{2018}{University of Pennsylvania}

\declaretheoremstyle[bodyfont=\normalfont]{normalbody}
\newtheorem{theorem}{Theorem}[section]
\newtheorem{lemma}[theorem]{Lemma}
\newtheorem*{claim}{Claim}
\newtheorem*{fact}{Fact}
\newtheorem{propdef}[theorem]{Proposition-Definition}
\newtheorem{corollary}[theorem]{Corollary}
\newtheorem{proposition}[theorem]{Proposition}
% \newtheorem{remark}[theorem]{Remark}
% \newtheorem{conjecture}[theorem]{Conjecture}

\newtheorem{definition}[theorem]{Definition}
% \newtheorem{example}[theorem]{Example}
\newtheorem{conjecture}[theorem]{Conjecture}
\newtheorem{xca}[theorem]{Exercise}
% \theoremstyle{remark}

% \theoremstyle{definition}
\theoremstyle{remark}
\declaretheorem{remark}
\declaretheorem[style=normalbody]{example}



% \numberwithin{equation}{section}

%\input{newcommands}
\begin{document}
\pagenumbering{roman}
% \frontmatter
% \title[short text for running head]{full title}
% \title{$f$-Stability and its Application in Differential Privacy}
\title{Gaussian Differential Privacy}
% \vspace{10cm}
% : A Primal View of $(\ep,\delta)$-Differential Privacy \ar{What does ``a primal view'' mean here?}}

%    Only \author and \address are required; other information is
%    optional.  Remove any unused author tags.

%    author one information
% \author[short version for running head]{name for top of paper}
\author{Jinshuo Dong\thanks{Graduate Group in Applied Mathematics and Computational Science, University of Pennsylvania. Email: \texttt{jinshuo@sas.upenn.edu}.} \and Aaron Roth\thanks{Department of Computer and Information Sciences, University of Pennsylvania. Email: \texttt{aaroth@cis.upenn.edu}.} \and Weijie J.~Su\thanks{Department of Statistics, the Wharton School, University of Pennsylvania. Email: \texttt{suw@wharton.upenn.edu}.}}
% \author{Jinshuo Dong, Aaron Roth, Weijie J.~Su}
% \address{Applie Mathematics and Computational Science, University of Pennsylvania}
% \curraddr{}
% \email{}
% \thanks{}

%    author two information
% \author{}
% \address{}
% \curraddr{}
% \email{}
% \thanks{}

%    \subjclass is required.
% \subjclass[2018]{Differential privacy, information theory}


%% do not use today.
\date{May 24, 2019}

% \dedicatory{}

\maketitle

\vspace{0.3cm}
\begin{abstract}
In the past decade, differential privacy has seen remarkable success as a rigorous and practical formalization of data privacy. This privacy definition and its divergence based relaxations, however, have several acknowledged weaknesses, either in handling composition of private algorithms or in analyzing important primitives like privacy amplification by subsampling. Inspired by the hypothesis testing formulation of privacy, this paper proposes a new relaxation of differential privacy, which we term ``$f$-differential privacy'' ($f$-DP). This notion of privacy has a number of appealing properties and, in particular, avoids difficulties associated with divergence based relaxations. First, $f$-DP faithfully preserves the hypothesis testing interpretation of differential privacy, thereby making the privacy
guarantees easily interpretable. In addition, $f$-DP allows for lossless reasoning about composition in an algebraic fashion. Moreover, we provide a powerful technique to import existing results proven for the original differential
privacy definition to $f$-DP and, as an application of this technique, obtain a simple and easy-to-interpret theorem of privacy amplification by subsampling for $f$-DP.


In addition to the above findings, we introduce a canonical single-parameter family of privacy notions within the $f$-DP class that is referred to as ``Gaussian differential privacy'' (GDP), defined based on hypothesis testing of two shifted Gaussian distributions. GDP is the focal privacy definition among the family of $f$-DP guarantees due to a central limit theorem for differential privacy that we prove. More precisely, the privacy guarantees of \emph{any} hypothesis testing based definition of privacy (including the original differential privacy definition) converges to GDP in the limit under composition. We also prove a Berry--Esseen style version of the central limit theorem, which gives a computationally inexpensive tool for tractably analyzing the exact composition of private algorithms.

Taken together, this collection of attractive properties render $f$-DP a mathematically coherent, analytically tractable, and versatile framework for private data analysis. Finally, we demonstrate the use of the tools we develop by giving an improved analysis of the privacy guarantees of noisy stochastic gradient descent.
\end{abstract}

\newpage
\tableofcontents

\newpage
\pagenumbering{arabic}
\setcounter{page}{1}
\input{intro}
% root = ../paper.tex
\section{$f$-Differential Privacy and Its Basic Properties}
\label{sec:fDP}


%\djs{Replace $\T$ notation.}
%\djs{Revise roadmap after we agree on the content of this section.}



In Section~\ref{sub:trade_off_function}, we give a formal definition of $f$-DP. Section~\ref{sub:gaussian_differential_privacy} introduces Gaussian differential privacy, a special case of $f$-DP. In Section~\ref{sec:conn-with-blackw}, we highlight some appealing properties of this new privacy notation from an information-theoretic perspective. Next, Section~\ref{sub:a_primal_dual_connection_with_} offers a profound connection between $f$-DP and $(\epsilon, \delta)$-DP. Finally, we discuss the group privacy properties of $f$-DP.

% Finally, we discuss connections between $f$-DP and Blackwell's theorem, $(\epsilon, \delta)$-DP, and its ``group privacy'' properties in %Sections~\ref{sec:conn-with-blackw}, \ref{sub:a_primal_dual_connection_with_}, and \ref{sub:group_privacy}, respectively.

% %%Then we discuss its properties and relation with other differential privacy notions such as $(\ep,\delta)$-DP and {\Renyi} DP.
% %%The central novelty of our theory is in the characterization of indistinguishablility: we switch from various divergences to \textit{trade-off functions}, which is
%introduced in \Cref{sub:trade_off_function}. \djs{Here or intro?}
%Before that, let's first clarify a few terms for readers not familiar with differential privacy.
% and then start from the motivation and definition of the central concepts of this paper in \cref{sub:trade_off_function}.

%%As a central idea, differential privacy requires that neighboring datasets should be almost indistinguishable.
%%Such a dataset is often organized as a $n\times d$ matrix, with rows corresponding to individuals. In this view,
%%a neighboring dataset of $\vec{x}$ is a matrix $\vec{x}'$ that differs $\vec{x}$ by one row. For example, $\vec{x}' = (x_1',\ldots, x_n)$ is a neighbor of $\vec{x}$ as %long as $x_1'\neq x_1$. What we just described is replace-one neighboring. There is also add-one/remove-one neighboring.


%%\djs{Changed to dataset. We might want to distinguish between collection and ordered collection. Now we claim a dataset is a collection, but our mathematical %representation is an ordered collection.}
%% \djs{Still not very consistent. I use $X^n$ in Sections 4 and 5 without explicitly saying $X=\R^d$.} Without loss of generality, a dataset $x$ can be thought of as a %collection of $n$ records, each consisting of a $d$-dimensional real-valued vector \ar{Any reason to make it a $d$-dimensional real valued vector? Why not have an %abstract data domain $\mathcal{X}$?}. We write:

Before moving on, we first establish several key pieces of notation from the differential privacy literature.
\begin{itemize}
\item  % (fold)\label{par:dataset}
{\bf Dataset.} A dataset $S$ is a collection of $n$ records, each corresponding to an individual. Formally, we write the dataset as $S = (x_1,\ldots, x_n)$, and an individual $x_i \in X$ for some abstract space $X$. Two datasets $S' = (x'_1,\ldots, x'_n)$ and $S$ are said to be \textit{neighbors} if they differ in exactly one record, that is, there exists an index $j$ such that $x_i = x_i'$ for all $i \ne j$ and $x_j \ne x_j'$.

%\wjs{say $Y$}\djs{Removed ``by appropriate permuting''. We do not explicit say our algorithms are invariant under the action of permutation group, and we don't have to.}

% Replace-$k$ neighboring corresponds to group privacy (see \Cref{sub:group_privacy}). In order to minimize such detail, it is helpful to make the following abstraction:
% % \begin{definition} \label{def:dataset}
% \begin{center}
% 	A \emph{dataset space} is an abstract set $X$ equipped with a symmetric binary relation $\sim$.
% \end{center}
% \djs{Putting it into a definition environment breaks the ``Database'' and ``Mechanism'' paragraph environments.}
% % \end{definition}
%%$X$ is a collection of possible instances of datasets. The neighboring relation is encoded in $\sim$. Note that a transitive $\sim$ typically does not lead to a meaningful privacy notion. \djs{More should be said on this in intro.}
% \ar{In fact, its not transitive in any interesting setting, right?}
% paragraph dataset (end)

\item {\bf Mechanism.}  A mechanism $M$ refers to a randomized algorithm that takes as input a dataset $S$ and releases some (randomized) statistics $M(S)$ of the dataset in some abstract space $Y$. For example, a mechanism can release the average salary of individuals in the dataset plus some random noise.

% (fold) \label{par:mechanism}
%\ar{Is this definition useful/correct? It doesn't have to output ``summary statistics'' -- the range can be anything. All we are saying is mechanism = algorithm. }
% . For some readers it may be helpful to think of a mechanism $M$ as a Markov kernel or conditional probability. When there is no ambiguity, the notation often intended for mappings, $M:X\to Y$, is abused to mean that $M$ is a mechanism/randomized algorithm whose input and output are elements of $X$ and $Y$ respectively.
% More precisely, when $M:X\to Y$ is fed with two neighboring datasets $x\sim x'$, the results $M(x)$ and $M(x')$ are distributions on $Y$. For example, $M:\R^{n\times d}\to\R^d$ may take the average of each column and add Gaussian noise to the average vector. $M(x)$ and $M(x')$ will be two $d$-dimensional Gaussian with different means in this case.
% % paragraph mechanism (end)

\end{itemize}

%%\subsection{Trade-off functions and Testing Regions} % (fold)

% root = ../paper.tex
\subsection{Trade-off Functions and $f$-DP}
\label{sub:trade_off_function}

%The essence of any existing differential privacy notion requires that it is fundamentally difficult to distinguish between two \textit{neighboring} datasets based on the releasing information of a private a mechanism $M$. As a consequence of this, differential privacy prevents from inferring any individual of the dataset using the releasing information, thus leading to the privacy of the mechanism. More precisely, for any two neighboring datasets $x$ and $x'$, $(\ep,\delta)$-DP formalizes the indistinguishability by requiring $M(x)$ and $M(x')$ be close in the sense of probability distribution, and the {\Renyi} DP uses the {\Renyi} divergence to measure the discrepancy between $M(x)$ and $M(x')$.
% To introduce $f$-DP, it is helpful to formalize this notion of 

All variants of differential privacy informally require that it be hard to \textit{distinguish} any pairs of neighboring datasets based on the information released by a private a mechanism $M$. From an attacker's perspective, it is natural to formalize this notion of  ``indistinguishability'' as a hypothesis testing problem for two neighboring datasets $S$ and $S'$:
\[
H_0: \text{the underlying dataset is } S \quad\text{ versus }\quad H_1: \text{the underlying dataset is } S'.
\]
The output of the mechanism $M$ serves as the basis for performing the hypothesis testing problem. Denote by $P$ and $Q$ the probability distributions of the mechanism applied to the two datasets, namely $M(S)$ and $M(S')$, respectively. The fundamental difficulty in distinguishing the two hypotheses is best delineated by the \textit{optimal} trade-off between the achievable type I and type II errors. More precisely, consider a rejection rule $0 \le \phi \le 1$, with type I and type II error rates defined as\footnote{A rejection rule takes as input the released results of the mechanism. We flip a coin and reject the null hypothesis with probability $\phi$.}
\[
\alpha_\phi = \E_{P}[\phi], \quad \beta_\phi = 1 - \E_{Q}[\phi],
\]
respectively. The two errors satisfy, for example, the constraint is well known to satisfy
% \djs{Removed the claim that there is a proof in the appendix.}
% (see its proof in \Cref{app:fDP}) \djs{Do we need a proof? Isnt' it well-known?}
\begin{equation}\label{eq:tv_norm}
\alpha_{\phi} + \beta_{\phi} \ge 1 - \TV(P, Q),
\end{equation}
where the total variation distance $\TV(P, Q)$ is the supremum of $|P(A) - Q(A)|$ over all measurable sets $A$. Instead of this rough constraint, we seek to characterize the fine-grained trade-off between the two errors. Explicitly, fixing the type I error at \textit{any} level, we consider the minimal achievable type II error. This motivates the following definition.
\begin{definition}[trade-off function] \label{def:trade-off}
For any two probability distributions $P$ and $Q$ on the same space, define the trade-off function $\F(P, Q): [0, 1] \rightarrow [0, 1]$ as
\[
\F(P, Q)(\alpha) = \inf \left\{ \beta_\phi: \alpha_\phi \leqslant \alpha \right\},
\]
where the infimum is taken over all (measurable) rejection rules.
\end{definition}
The trade-off function serves as a clear-cut boundary of the achievable and unachievable regions of type I and type II errors, rendering itself the \textit{complete} characterization of the fundamental difficulty in testing between the two hypotheses. In particular, the greater this function is, the harder it is to distinguish the two distributions. For completeness, we remark that the minimal $\beta_{\phi}$ can be achieved by the likelihood ratio test---a fundamental result known as the Neyman--Pearson lemma, which we state in the appendix as \Cref{thm:NPlemma}.




A function is called a trade-off function if it is equal to $T(P, Q)$ for some distributions $P$ and $Q$. Below we give a necessary and sufficient condition for $f$ to be a trade-off function. This characterization reveals, for example, that $\max\{f, g\}$ is a trade-off function if both $f$ and $g$ are trade-off functions.
\begin{restatable}{proposition}{tradeoffthm}\label{prop:trade-off}
A function $f: [0, 1] \rightarrow [0, 1]$ is a trade-off function if and only if $f$ is convex, continuous\footnote{Convexity itself implies continuity in $(0,1)$ for $f$. In addition, $f(\alpha) \ge 0$ and $f(\alpha) \le 1-\alpha$ implies continuity at 1. Hence, the continuity condition only matters at $x = 0$.}, non-increasing, and $f(x) \leqslant 1-x$ for $x \in [0,1]$.
\end{restatable}
%Throughout the paper, $f$ is used to denote a trade-off function.

%%In the case where $P = Q$, we use $\Id$ to denote\footnote{Identity function $\alpha\mapsto \alpha$ is the identity element for function composition. In contrast, our $\Id$ has a different expression, and works as the identity element in a different algebraic structure. See \Cref{sub:composition_theorem}.} the trade-off function $\F(P , P)(\alpha) = 1 - \alpha$.

% Optimal tests achieve minimal type II error $\beta$ for a fixed level of type I error $\alpha$. This yields a correspondence $\alpha\mapsto \beta(\alpha)$ that characterizes the ``trade-off'' between type I and type II errors. We will call this the trade-off function of $P_0$ and $P_1$, and denote it by $\F(P_0,P_1)$, taking the initial of ``trade-off''.
%\bigskip
% The corresponding type I and type II error have expressions:
% \[\beta_\phi = 1-\E_{P_1}[\phi],\quad\alpha_\phi = \E_{P_0}[\phi].\]
% With these notations,
% % That is, $\F(P_0,P_1)$ takes in a level $\alpha\in[0,1]$ and outputs the corresponding minimal type II error.
% The infimum is taken over all rejection rules, i.e. over all measurable functions $\phi:\Omega\to[0,1]$. Neyman-Pearson lemma says for any $\alpha\in[0,1]$, there exists a likelihood ratio test that achieves type I error $\alpha$ and type II error $\F(P_0,P_1)(\alpha)$, so the infimum can be written as minimum.
%Formally, if the sample space of $P_0$ and $P_1$ is $\Omega$, then statistical tests, or equivalently rejection rules, have the form: $\phi:\Omega\to[0, 1]$. Suppose the outcome is $\omega\in\Omega$. The interpretation is
%Let $(X,\sim)$ be a dataset space and $Y$ an arbitrary set.

%This function does characterize the indistinguishability of $P_0$ and $P_1$:
%Roughly speaking,


% %To better appreciate this trade-off, consider two distributions $\widetilde P$ and $\widetilde Q$ that have the same trade-off function as $P$ and $Q$, that is, $\F(\widetilde P, \widetilde Q)(\alpha) = \F(P, Q)(\alpha)$ for all $0 \le \alpha \le 1$. As is clear, distinguishing between $P$ and $Q$ is \textit{equally} difficult as distinguishing between $\widetilde P$ and $\widetilde Q$. As such, the trade-off function is \textit{complete} for characterizing the hardness of testing two distributions from an information-theoretic perspective. \ar{I don't understand this claim. Does
%   it mean something formal? Otherwise I suggest removing.}\wjs{I will do it soon.}

Now, we propose a new generalization of differential privacy built on top of trade-off functions. Below, we write $g \ge f$ for two functions defined on $[0, 1]$ if $g(x) \ge f(x)$ for all $0 \le x \le 1$, and we abuse notation by identifying $M(S)$ and $M(S')$ with their corresponding probability distributions. Note that if $T(P,Q) \ge T(\widetilde P, \widetilde Q)$, then in a very strong sense, $P$ and $Q$ are harder to distinguish than $\widetilde P$ and $\widetilde Q$ at \textit{any} level of type I error.

%\djs{A better sentence for this?} In this spirit, our $f$-DP imposes a pointwise lower bound $f$ on $T\big(M(x),M(x')\big)$.
%%$\geqslant$ is abused for pointwise relation here, and will be used this way through out this paper.
%%A meaningful $f$ should have range in $[0,1]$, as the range of $T\big(M(x),M(x')\big)$ is always in $[0,1]$.
% It's now an object on its own right, no longer associated with the pair $(P_0,P_1)$, which is not unique. For example, $N(0,1)$ and $N(\mu,1)$ obviously has the same trade-off function as $N(0,1)$ and $N(-\mu,1)$.
\begin{definition}[$f$-differential privacy] \label{def:kstable}
Let $f$ be a trade-off function. A mechanism $M$ is said to be $f$-differentially private if
\[
T\big(M(S), M(S')\big) \ge f
\]
for all neighboring datasets $S$ and $S'$.
\end{definition}



A graphical illustration of this definition is shown in Figure~\ref{fig:intro_def}. Letting $P$ and $Q$ be the distributions such that $f = \F(P, Q)$, this privacy definition amounts to saying that a mechanism is $f$-DP if distinguishing any two neighboring datasets based on the released information is at least as difficult as distinguishing $P$ and $Q$ based on a single draw. In contrast to existing definitions of differential privacy, our new definition is parameterized by a function, as opposed to several real valued parameters (e.g.~$\epsilon$ and $\delta$). This functional perspective offers a complete characterization of ``privacy'', thereby avoiding the pitfall of summarizing statistical information too early. This fact is crucial to the development of a composition theorem for $f$-DP in \Cref{sec:composition-theorems}. Although this completeness comes at the cost of increased complexity, as we will see in Section~\ref{sub:gaussian_differential_privacy}, a simple family of trade-off functions can often closely capture privacy loss in many scenarios.

\begin{figure}[!htp]
\centering
\includegraphics[width=.6\linewidth]{./figures/nonfDP.pdf}
\caption{Three different examples of $T\big(M(S),M(S')\big)$. Only the dashed line corresponds to a trade-off function satisfying $f$-DP.}
\label{fig:intro_def}
\end{figure}

%It also provides flexibility the choice of the function $f$. This flexibility is particularly useful in the case where the indistinguishability between $M(x)$ and $M(x')$ can be well approximated by a small family of trade-off functions (see Section~\ref{sub:gaussian_differential_privacy}). In relation to the literature, we remark that \cite{KOV} uses testing regions instead of trade-off functions as a measure of the hardness of hypothesis testing (see \Cref{app:fDP} for the equivalence between the two perspectives). The perspective of trade-off functions brings much more convenience in the development of the composition
%theory of $f$-DP in Section~\ref{sec:composition_theorem}. \wjs{Jinshuo: could you please check?}

% Denote by $\tv(f)$ the total variation of the two distributions that realize $f$. From \eqref{eq:tv_norm}, we get the following implication of $f$-DP.
% \begin{proposition}\label{prop:dTV}
% Let $M$ be an $f$-DP mechanism. Then, the type I and type II errors for the hypothesis testing problem of distinguishing $M(S)$ from $M(S')$ for neighboring $S,S'$ sum to at least $1 - \tv(f)$.
% \end{proposition}

%As is clear, $M$ is $f$-DP implies that $M$ is $g$-DP for a trade-off function $g$ that satisfies $g \preceq f$.
% However, it is much more convenient to develop the theory in terms of function, since at various stages we want to do differential and integral calculus, take (function) composition, etc.
%\wjs{in hypotheis testing, null and non-null are not symmetric.}
%A leftover question in the definition of $f$-DP is, the symmetry of the neighboring relation is not revealed in $f$. For example, if $M$ is $f$-DP, then we know both of the following are true:
%It is fair to say that any reasonable privacy notion must have this ``robust under post-processing'' property, i.e., a mechanism only becomes more private when processed \textbf{independent of individuals' data}. Intuitively, any information extractable from the processed output must be extractable from the original output. This is a natural requirement, and it (almost) automatically holds for every version of differential privacy ever appeared, including our $f$-DP.
% Since previous differential privacy  notions are all defined in terms of some divergence, the corresponding post-processing property follows naturally from the corresponding data processing inequality. \djs{What do I want to say...}
%%Suppose we have a mechanism $M$ from a dataset space $X$ to some $Y$. When we apply a post-processing $\Proc :Y\to Z$ we get a new mechanism $\Proc \circ M:X\to Z$. Note that $\Proc $ only accepts a $y\in Y$ but not the dataset as input. Rigorously speaking, $M$ and $\Proc $ should be interpreted as Markov kernels and $\circ$ is Markov kernel composition.
%%ntuitively, the processed output $\Proc \circ M(x)$ releases no more information than $M(x)$ about the dataset $x$ and, therefore, is more private.


%Next, we present some simple properties of this new privacy definition.
%\noindent{\bf Symmetric trade-off functions.}

%Recognizing that the two hypotheses regarding the neighboring $S$ and $S'$, the trade-off function $f$ had better reflect a certain symmetry on its own. To put this point into perspective, note that the definition of $f$-DP also requires

Naturally, the definition of $f$-DP is symmetric in the same sense as the neighboring relationship, which by definition is symmetric. Observe that this privacy notion also requires
\[
T\big(M(S'), M(S)\big) \ge f
\]
for any neighboring pair $S, S'$. Therefore, it is desirable to restrict our attention to ``symmetric'' trade-off functions. Proposition~\ref{prop:symmetry} shows that this restriction does not lead to any loss of generality.
\begin{restatable}{proposition}{symmrep}\label{prop:symmetry}
Let a mechanism $M$ be $f$-DP. Then, $M$ is $\Sf$-DP with $\Sf = \max\{f, f^{-1}\}$, where the inverse function is defined as\footnote{
\Cref{eq:inver_f} is the standard definition of the left-continuous inverse of a decreasing function. When $f$ is strictly decreasing and $f(0)=1$ and hence bijective as a mapping, \eqref{eq:inver_f} corresponds to the inverse function in the ordinary sense, i.e. $f(f^{-1}(x)) = f^{-1}(f(x)) = x$. However, this is not true in general.}
\begin{equation}\label{eq:inver_f}
f^{-1}(\alpha):=\inf\{t \in [0,1]:f(t) \leqslant \alpha\}
\end{equation}
for $\alpha \in[0,1]$.
\end{restatable}
%\wjs{Now $f$ is not necessarily trade-off, is $f^{-1}$ still well-defined? Does $(f^{-1})^{-1} = f$ still hold? Does this proposition remain true?}
%, which is the left continuous inverse of $f$.
%. In addition, $\tilde{f}\in\T$ and $\tilde{f}^{-1} = \tilde{f}$.%justifies this intuition.
% A trade-off function $f\in\T$ is not necessarily 1-1 onto, so it takes some care to properly define $f^{-1}$. It turns out the proper language is the left continuous inverse. Namely,
% Note that in general $f^{-1}\big(f(x)\big) = x$ is not true.
% \djs{$f^{-1}$ in $\T$?}
% The class of symmetric trade-off function is defined to be
% $$\T^S:=\{f\in\T: f^{-1} = f\}.$$

Writing $f = \F(P,Q)$, we can express the inverse as $f^{-1} = \F(Q, P)$, which therefore is also a trade-off function. As a consequence of this, $\Sf$ continues to be a trade-off function by making use of \Cref{prop:trade-off} and, moreover, is \textit{symmetric} in the sense that 
\[
\Sf = (\Sf)^{-1}.
\] 
Importantly, this symmetrization gives a tighter bound in the privacy definition since $\Sf \geqslant f$. In the remainder of the paper, therefore, trade-off functions will always be assumed to be symmetric unless otherwise specified. We prove \Cref{{prop:symmetry}} in \Cref{app:fDP}.

% In short, this proposition reveals that $f$ can always be assumed to be symmetric.%\\[0.2em]

%\noindent{\bf $f$-DP is an extension of $(\epsilon, \delta)$-DP.}

We conclude this subsection by showing that $f$-DP is a generalization of $(\ep,\delta)$-DP. This foreshadows a deeper connection between $f$-DP and $(\epsilon, \delta)$-DP that will be discussed in Section~\ref{sub:a_primal_dual_connection_with_}. Denote
\begin{equation}\label{eq:fed}
f_{\ep, \delta}(\alpha) = \max\left\{ 0,1 - \delta - \e^\ep \alpha, \e^{-\ep}(1-\delta-\alpha) \right\}
\end{equation}
for $0 \le \alpha \le 1$, which is a trade-off function. \Cref{fig:DPvsGDP} shows the graph of this function and its evident symmetry. The following result is adapted from \cite{wasserman_zhou}.
%It must be truncated at 0 to obtain the residency in $\T$.
\begin{proposition}[\cite{wasserman_zhou}] \label{thm:privacy_testing}
A mechanism $M$ is $(\ep, \delta)$-DP if and only if $M$ is $f_{\ep, \delta}$-DP.
\end{proposition}



\begin{figure}[!htp]
\centering
  \includegraphics[width=.75\linewidth]
{./figures/DPvsGDP.pdf}
  \captionof{figure}{Left: $f_{\ep,\delta}$ is a piecewise linear function and is symmetric with respect to the line $y = x$. It has (nontrivial) slopes $-\e^{\pm\ep}$ and intercepts $1-\delta$. Right: Trade-off functions of  unit-variance Gaussian distributions with different means. The case of $\mu=0.5$ is reasonably private, $\mu=1$ is borderline private, and $\mu=3$ is basically non-private: an adversary can control type I and type II errors simultaneously at only 0.07. In the case of $\mu=6$ (almost coincides with the axes), the two errors both can be as small as 0.001.}
  \label{fig:DPvsGDP}
\end{figure}



%For this reason, we will always assume a symmetric $f$ in the definition of $f$-DP.
%this new definition provides privacy measure is accurate and faithful hardly loses any characterization gives a more comphrenensive of. Meanhwile, GDP does not lose
%Apparently we should add more assumptions on $f$ to exclude pathological functions.
%If $f$ satisfies them, we will simply say $f$ is a trade-off function. Now trade-off function is no longer just a companion of distributions but an object in its own right. The collection of trade-off functions is denoted by $\T$.
%By the above proposition, $T\big(M(x),M(x')\big)$ is always convex and non-increasing, so focusing on lower bounds of the same type is not too restricted.
% \djs{It would be more convincing if the following were true: if a trade-off function $g$ has a general lower bound $f$, then it has a trade-off lower bound $\tilde{f}$ with $\tilde{f}\geqslant f$.}
% Although a theory for general $f$ would be preferable, the nature of the class $\T$ turns out to be important in subsampling and composition. The extra interpretability is another bonus.
%The symmetry of the neighboring relation ``$x\sim x'$ if and only if $x'\sim x$'' should also be revealed in $f$. We put off the detail until \Cref{ssub:symmetric_trade_off_functions}.
%The proof of \Cref{thm:trade-off} is not as important as its content,

% \begin{definition}
% 	A mechanism $M:X\to Y$ is $(\ep, \delta)$-DP if for any $x\sim x'$ and measurable $E\subseteq Y$,
% 	$$P[M(x)\in E] \leqslant e^\ep P[M(x')\in E] + \delta.$$
% \end{definition}
% Now the question is, what kind of distributions generates the testing region that corresponds to $(\ep, \delta)$-indistinguishability? Consider the simple measurable space $\{a,b,c,d\}$ and two probability measures $P_U$ and $P_V$ given in Table 1 \djs{ref not work?}:
% \begin{table}
% \centering
% \label{tab:PUPV}
% \begin{tabular}{|c|c|c|}
% \hline
% outcome & $P_U$ & $P_V$ \\\hline
% $a$ & $(1-\delta)\cdot\epover$ & $(1-\delta)\cdot\oneover$ \\\hline
% $b$ & $(1-\delta)\cdot\oneover$ & $(1-\delta)\cdot\epover$ \\\hline
% $c$ & $\delta$ & $0$ \\\hline
% $d$ & $0$ & $\delta$\\\hline
% \end{tabular}
% \caption{$P_U$ and $P_V$}
% \end{table}
% \begin{lemma} \label{lem:epspower}
% 	$K(P_U,P_V) = K(P_V,P_U)$ is the convex hull of the following five points: $(0,1), (1,0), (0,1-\delta), (1-\delta,0)$ and $ (,)$. It is (in general) a pentagon framed by the triangle $\triangle$ and two lines with slopes $-e^\ep$ and $-e^{-\ep}$.
% \end{lemma}
% \djs{Change when the plot is ready.}
% $(\ep, \delta)$ fully determines $(P_U,P_V)$, so we will denote $K(P_U,P_V)$ by $K_{\ep, \delta}$.


% \begin{proof}
% %\djs{Necessary?}
% \end{proof}
% subsection indistinguishability_based_on_testing_regions (end)




%%% symmetry as a subsubsection
%\input{fDP/symmetry}

%%% Local Variables:
%%% mode: latex
%%% TeX-master: "../paper"
%%% End:

\input{fDP/GDP}
\input{fDP/post}
\input{fDP/dual}
\input{fDP/group}



%\input{fDP/group}
% subsection a_primal_dual_connection_with_ (end)
%%% Local Variables:
%%% mode: latex
%%% TeX-master: "../paper"
%%% End:

\input{CLT/composition}
% root = paper.tex
\section{Amplifying Privacy by Subsampling}
\label{sec:subsampling}

\newcommand{\Sample}{\mathtt{Sample}}
\newcommand{\tm}{\widetilde{M}}
\newcommand{\Sn}{S\cup\{n\}}
\newcommand{\xv}{S}
% \subsection{Subsampling} % (fold)
%

Subsampling is often used prior to a private mechanism $M$ as a way to \textit{amplify} privacy guarantees. Specifically, we can construct a smaller dataset $\tilde S$ by flipping a fair coin for each individual in the original dataset $S$ to decide whether the individual is included in $\tilde{S}$. This subsampling scheme roughly shrinks the dataset by half and, therefore, we would expect that the induced mechanism applied to $\tilde S$ is about twice as private as the original mechanism $M$. Intuitively speaking, this privacy amplification is due to the fact that every individual enjoys perfect privacy if the individual is not included in the resulting dataset $\tilde S$, which happens with probability 50\%.

% \wjs{{\Renyi} Differential Privacy, {\Renyi} differential privacy, {\Renyi} DP. Should be consistent.}
The claim above was first formalized in \cite{KLNRS} for $(\ep,\delta)$-DP. Such a privacy amplification property is, unfortunately, no longer true for the most natural previous relaxations of differential privacy aimed at recovering precise compositions (like concentrated differential privacy (CDP) \cite{concentrated,concentrated2}). Further modifications such as truncated CDP \cite{tcdp} have been introduced primarily to remedy this deficiency of CDP---but at the cost of extra complexity in the definition. Other relaxations like {\Renyi} differential privacy \cite{renyi} can be shown to satisfy a form of privacy amplification by subsampling, but both the analysis and the statement are complex \cite{wang2018subsampled}.

% \ar{Here is the old paragraph. What part of it is inaccurate? \Renyi-DP and tCDP -do- have amplification theorems, right? So your new paragraph which says that ``all'' composition friendly relaxations are not compatible with the amplification property seems misleading.. }

In this section, we show that these obstacles can be overcome by our hypothesis testing based relaxation of differential privacy. Explicitly, our main result is a simple, general, and easy-to-interpret subsampling theorem for $f$-DP. Somewhat surprisingly, our theorem significantly improves on the classical subsampling theorem for privacy amplification in the $(\ep,\delta)$-DP framework \cite{jon}. Note that this classical theorem continues to use $(\epsilon,\delta)$-DP to characterize the subsampled mechanism. However, $(\epsilon,\delta)$-DP is simply not expressive enough to capture the amplification of privacy.

% (by no longer expressing the resulting guarantees using $(\epsilon,\delta)$-differential privacy).
% in some parameter regimes.



\subsection{A Subsampling Theorem} % (fold)
\label{sec:subsample_theorems}

% section subsample_theorems (end)
%\djs{Determine the notation in earlier sections first.} Suppose the datum of each individual lies in an abstract domain $X$ (e.g. $X=\R^d$ or $X=\{0,1\}^d$). When there are $n$ individual, our dataset space is $X^n$. A particular dataset is denoted by $\xv = (x_1,\ldots, x_n)$. By individual we refer to an index in $\{1,2,...,n\}$. $x_i$ is the datum that belongs to individual $i$. Hence mechanisms of interest in this section are random mappings from $X^n$ to some abstract range $Y$. \ar{This is how we should define datasets/mechanisms initially.}
	% \item For each of the $n$ individuals in $\xv$, include the individual's datum in the subsample independently with probability $p$. The resulting subsample is denoted by $\Sample(\xv,p)$;.
 %It does not matter because our theory is based on a general reduction from subsampling theorems for $(\ep,\delta)$-DP to those for $f$-DP. Whatever is present for $(\ep,\delta)$-DP will be present for $f$-DP. \djs{Added this based on conversation with Aaron.}
%	\item There are $n \choose m$ different subsets of the $n$ individuals with size $m$. Uniformly pick one from these $n \choose m$ subsets. The pick is (abusively) referred to as $\Sample(\xv,m)$.
% Since $\Sample(\xv,p)$ can have uncertain size, in this case $M$ has to operate on $X^{\leqslant n}:=\bigcup_{i=0}^n X^i$. The neighboring relation is pieced together on each $X^i$: two datasets are neighbors if they differ by exactly one individual. In particular, if they are of different size then they are not neighbors.
% Suppose we are given that $M$ is $f$-DP. After subsampling we get the mechanism $\tm(\xv)=M\big(\Sample(\xv,p)\big)$ (or $\tm(\xv)=M\big(\Sample(\xv,m)\big)$ depending on the sampling scheme). It is convenient to hide the input $\xv$ by abusing the deterministic mapping composition notation for randomized mappings and write $\tilde{M} = M\circ\Sample(p)$.

Given an integer $1 \le m \le n$ and a dataset $S$ of $n$ individuals, let $\Sample_m(\xv)$ be a subset of $S$ that is chosen uniformly at random among all the $m$-sized subsets of $\xv$. For a mechanism $M$ defined on $X^m$, we call $M\big(\Sample_m(\xv)\big)$ the subsampled mechanism, which takes as input an $n$-sized dataset. Formally, we use $M \circ \Sample_m$ to denote this subsampled mechanism. To clear up any confusion, note that intermediate result $\Sample_m(\xv)$ is not released and, in particular, this is different from the composition in \Cref{sec:composition-theorems}.



In brief, our main theorem shows that the privacy bound of the subsampled mechanism in the $f$-DP framework is given by an operator acting on trade-off functions. To introduce this operator, write the convex combination $f_p:= pf + (1-p)\Id$ for $0 \le p \le 1$, where $\Id(x) = 1-x$. Note that the trade-off function $f_p$ is asymmetric in general.
\begin{definition} \label{def:Cp}
For any $0 \le p \le 1$, define the operator $C_p$ acting on trade-off functions as
\[C_p(f) := \min\{f_p, f_p^{-1}\}^{**}.\]
We call $C_p$ the $p$-sampling operator.
\end{definition}
Above, the inverse $f_p^{-1}$ is defined in \eqref{eq:inver_f}. The biconjugate $\min\{f_p, f_p^{-1}\}^{**}$ is derived by applying the conjugate as defined in \eqref{eq:conjugate} twice to $\min\{f_p, f_p^{-1}\}$. For the moment, take for granted the fact that $C_p(f)$ is a symmetric trade-off function.

% Basic properties of convex conjugates show that $C_p$ maps $f$ to a symmetric trade-off function:
% \begin{lemma}
% As a function defined on $[0, 1]$, $C_p(f)$ is a symmetric trade-off function provided that $f$ is a trade-off function.
% \end{lemma}
% % and is, therefore, convex. Consequently, $C_p(f)$ is a (symmetric) trade-off function.

Now, we present the main theorem of this section.

\begin{theorem}\label{thm:subsample}
If $M$ is $f$-DP on $X^m$, then the subsampled mechanism $M\circ\Sample_m$ is $C_p(f)$-DP on $X^n$, where the sampling ratio $p=\frac{m}{n}$.
\end{theorem}

%This privacy guarantee of the subsampled mechanism is tight in the sense that the theorem does not hold if $C_p(f)$ is replaced by any larger trade-off function. \wjs{This is true?}

Appreciating this theorem calls for a better understanding of the operator $C_p$. In effect, $C_p$ performs a two-step transformation: symmetrization (taking the minimum of $f_p$ and its inverse $f_p^{-1}$) and convexification (taking the largest convex lower envelope of $\min\{f_p, f_p^{-1}\}$). The convexification step is seen from convex analysis that the biconjugate $h^{**}$ of any function $h$ is the greatest convex lower bound of $h$. As such, $C_p(f)$ is convex and, with a bit more analysis, \Cref{prop:trade-off} ensures that $C_p(f)$ is indeed a trade-off function. As an aside, $C_p(f) \le \min\{f_p, f_p^{-1}\} \le f_p$. See \Cref{fig:subsample} for a graphical illustration.


%$= \mathrm{Conv}\big(\min\{f_p,f_p^{-1}\}\big) $
% Now, we present the main theorem of this section, which quantifies the privacy bound of the subsampled mechanism in the $f$-DP framework. In effect, subsampling
% In our $f$-DP framework, privacy is quantified by trade-off functions, so in our theorems, amplification incurred by subsampling with ratio $p$ is described by an operation $C_p$ which acts on trade-off functions.
%$C_p$ roughly does the following to a symmetric trade-off function $f$
%The formal definition of $C_p$ and discussion is deferred until after the main theorem of this section.
%Formalizing the description above, $C_p$ can be rigorously defined as follows:
% Conv is just a semantically more interpretable notation for it.
%understand how privacy is amplified by subsampling.
%lead to a qualitative understanding of privacy
%are we provide several easy to verify facts that help the reader better understand the qualitative behavior of this operator.


Next, the following facts concerning the $p$-sampling operator qualitatively illustrate this privacy amplification phenomenon.
\begin{enumerate}
\item If $0\leqslant p\leqslant q\leqslant 1$ and $f$ is symmetric, we have $f=C_1(f)\leqslant C_q(f)\leqslant C_p(f)\leqslant C_0(f)= \Id$. That is, as the sampling ratio declines from 1 to 0, the privacy guarantee interpolates monotonically between the original $f$ and the perfect privacy guarantee $\Id$. This monotonicity follows from the fact that $g \ge h$ is equivalent to $g^{-1} \ge h^{-1}$ for any trade-off functions $g$ and $h$.

\item
If two trade-off functions $f$ and $g$ satisfy $f \ge g$, then $C_p(f) \ge C_p(g)$. This means that if a mechanism is more private than the other, using the same sampling ratio, the subsampled mechanism of the former remains more private than that of the latter, at least in terms of lower bounds.

\item For any $0 \le p \le 1$, $C_p(\Id) = \Id$. That is, perfect privacy remains perfect privacy with subsampling.
\end{enumerate}

Explicitly, we provide a formula to calculate $C_p(f)$ for a symmetric trade-off function $f$. Letting $x^*$ be the unique fixed point of $f$, that is $f(x^*)= x^*$, we have
\begin{equation}\label{eq:Cp_expression}
		C_p(f)(x)=
			\left\{
			\begin{array}{ll}
			f_p(x),&x\in[0,x^*] \\
			x^*+f_p(x^*)-x, & x\in[x^*,f_p(x^*)]\\
			f_p^{-1}(x), &x\in[f_p(x^*),1].
			\end{array}
			\right.
	\end{equation}
This expression is almost self-evident from the left panel of \Cref{fig:subsample}. Nevertheless, a proof of this formula is given in \Cref{app:property}. This formula, together with \Cref{thm:subsample}, allows us to get a closed-form characterization of the privacy amplification for $(\epsilon, \delta)$-DP.
\begin{corollary}\label{cor:ep_d}
If $M$ is $(\epsilon, \delta)$-DP on $X^m$, then the subsampled mechanism $M\circ\Sample_m$ is $C_p(f_{\epsilon, \delta})$-DP on $X^n$, where
\begin{equation}\label{eq:ep_d_sump_for}
C_p(f_{\epsilon, \delta})(\alpha) = \max\left\{f_{\epsilon', \delta'}(\alpha), 1 - p\delta - p \, \frac{\e^{\ep} - 1}{\e^{\ep} + 1} - \alpha \right\}.
\end{equation}
Above, $\epsilon'= \log(1-p + p \e^\ep), \delta' = p\delta$, and $p = \frac{m}{n}$.
\end{corollary}

% Further details can be found in \Cref{app:property}.
% and provide the proof of expression \eqref{eq:Cp_expression} in \Cref{app:property}.
% where $x^*$ is the unique fixed point of $f$, i.e. $f(x^*)= x^*$.
% \end{proposition}
% The second break point $px^*+(1-p)(1-x^*)$ is equal to $f_p(x^*)$.
% \begin{proposition}\label{prop:Cp_expression}
% For a symmetric trade-off function $f$ and $p\in[0,1]$, $C_p$ has the following expression:





\begin{figure}[!htp]
\centering
  \includegraphics[width=.7\linewidth]
{./figures/subsample_main.pdf}
  \captionof{figure}{The action of $C_p$. Left panel: $f=G_{1.8}, p=0.35$. Right panel: $\ep = 3,\delta=0.1,p=0.2$. The subsampling theorem \ref{thm:subsample} results in a significantly tighter trade-off function compared to the classical theorem for $(\ep,\delta)$-DP.}
% Even if one wants to express the final guarantee of a composition of subsampled mechanisms using $(\ep,\delta)$-DP, expressing the guarantees of intermediate steps of a composition in this way can be unavoidably wasteful.
  \label{fig:subsample}
\end{figure}

% If we use $f_p$ to denote the convex combination $pf+(1-p)\Id$, then the expression for $C_p(f)$ is
% Let's use $f_p$ to denote the convex combination $pf+(1-p)\Id$ temporarily in this section. Not to confuse it with other ``trade-off function with subscript'' notation used in this paper. \djs{Put it somewhere in the beginning of the paper.} Thus $C_p(f)$ is the maximum convex function that is smaller than both $f_p$ and $f_p^{-1}$.
% The tightness of these theorems are less credible since there is no statement in terms of indistinguishability. Unlike composition theorems and group privacy, the universal quantifiers are necessary, even in the statement.
% $\tilde{M} = M\circ\Sample(p)$ is $(\ep',\delta')$-DP with $\ep' = \log(1-p + pe^\ep), \delta' = p\delta$.

For comparison, we now present the existing bound on the privacy amplification by subsampling for $(\ep,\delta)$-DP. To be self-contained, \Cref{app:property} gives a proof of this result, which primarily follows \cite{jon} .

\begin{restatable}[\cite{jon}]{lemma}{DPsubsamplerep}\label{lem:DPsubsample}
If $M$ is $(\ep,\delta)$-DP, then $M\circ\Sample_m$ is $(\ep',\delta')$-DP with $\ep'$ and $\delta'$ defined in \Cref{cor:ep_d}.
%Recall the sampling ratio $p=\frac{m}{n}$.
\end{restatable}

Using the language of the $f$-DP framework, \Cref{lem:DPsubsample} states that $M\circ\Sample_m$ is $f_{\epsilon',\delta'}$-DP. \Cref{cor:ep_d} improves on \Cref{lem:DPsubsample} because, as is clear from \eqref{eq:ep_d_sump_for},
\[
C_p(f_{\epsilon, \delta}) \ge f_{\epsilon',\delta'}.
\]
The right panel of \Cref{fig:subsample} illustrates \Cref{lem:DPsubsample} and our \Cref{cor:ep_d} for $\epsilon = 3, \delta = 0.1$, and $p = 0.2$. In effect, the improvement is captured by the shaded triangle enclosed by $C_p(f_{\epsilon, \delta})$ and $f_{\epsilon',\delta'}$, revealing that the minimal sum of type I and type II errors in distinguishing two neighboring datasets with subsampling can be significantly lower than the prediction of \Cref{lem:DPsubsample}. This gain is only made possible by the flexibility of trade-off functions in the sense that $C_p(f_{\epsilon, \delta})$ \textit{cannot} be expressed within the $(\epsilon, \delta)$-DP framework. The unavoidable loss in the $(\epsilon,\delta)$-DP representation of the subsampled mechanism is compounded when analyzing the composition of many private mechanisms. 

% it is   there is a large gain in terms of the guarantee on the minimal (type I $+$ type II) error that any adversary can achieve in distinguishing two pairs of neighboring datasets.
%The proof of \Cref{lem:DPsubsample} is non-trivial. , which primarily follows Ullman's lecture notes \cite{jon}. \wjs{\cite{KLNRS} also did this?}

In the next subsection, we prove \Cref{thm:subsample} by making use of \Cref{lem:DPsubsample}. Its proof implies that \Cref{thm:subsample} holds for any subsampling scheme for which \Cref{lem:DPsubsample} is true. In particular, it holds for the subsampling scheme described at the beginning of this section, that is, independent coin flips for every data item. 

%\djs{Changed to this.}
% Before moving on to the proof of \Cref{thm:subsample} in \Cref{sub:proof_of_subsample_theorems}, we briefly remark that this theorem can be extended to the subsampling scheme described at the beginning of this section, that is, independent coin flips for every data item. See \Cref{app:property} for more discussion.

% Note that this scheme differs from the one informally described at the beginning of this section, since it results in datasets of size \emph{exactly} $p\cdot n$, rather than independently sampling each individual with probability $p$. But a similar theorem (with essentially the same proof) holds for the independent sampling scheme. This is because our proof gives a general way to lift known subsampling theorems for $(\ep,\delta)$-DP to theorems for $f$-DP, and amplification theorems for both sampling schemes are known for $(\ep,\delta)$-DP.
	% From the above discussion, the new linear envelopes are averages of the old ones with the diagonal $\Id$. These new envelopes correspond to $f_p = pf+(1-p)\Id$. To see this, it suffices to


\subsection{Proof of the Subsampling Theorem} % (fold)
\label{sub:proof_of_subsample_theorems}
The proof strategy is as follows. First, we convert the $f$-DP guarantee into an infinite collection of $(\ep,\delta)$-DP guarantees by taking a dual perspective that is enabled by \Cref{prop:ftoDP}. Next, by applying the classical subsampling theorem (that is, \Cref{lem:DPsubsample}) to these $(\ep,\delta)$-DP guarantees, we conclude that the subsampled mechanism satisfies a new infinite collection of $(\ep,\delta)$-DP guarantees. Finally, \Cref{prop:DPtof} allows us to convert these new privacy guarantees back into an $\tilde{f}$-DP guarantee, where $\tilde f$ can be shown to coincide with $C_p(f)$.

%\wjs{For simplicity, all trade-off functions in this section are assumed to be smooth. All the statements carry over to general trade-off functions with a subdifferential argument.}
% Executing the simple plan requires a slightly more advanced version of primal-dual conversion than what we have seen in \Cref{sub:a_primal_dual_connection_with_}.
%%	We know $M$ is $f$-DP. Converting to the dual description by

\begin{proof}[Proof of \Cref{thm:subsample}]
% Without loss of generality, assume $f$ is symmetric. Otherwise, the symmetrization trick from \Cref{prop:symmetry} can improve the bound $C_p(f)$ in the present theorem to $C_p(\max\{f, f^{-1}\})$. Note that $C_p(\max\{f, f^{-1}\}) \ge C_p(f)$ since $\max\{f, f^{-1}\} \ge f$.
Provided that $M$ is $f$-DP, from \Cref{prop:ftoDP} it follows that $M$ is $\big(\ep,\delta(\ep)\big)$-DP with $\delta(\ep) = 1+f^*(- \e^{\ep})$ for all $\ep\geqslant 0$. Making use of \Cref{lem:DPsubsample}, the subsampled mechanism $M \circ \Sample_m$ satisfies the following collection of $(\ep',\delta')$-DP guarantees for all $\ep\geqslant 0$:
\[
\ep' =\log(1 - p + p\e^\ep),\quad \delta' =p\big(1+f^*(-\e^{\ep})\big).
\]
Eliminating the variable $\epsilon$ from the two parametric equations above, we can relate $\ep'$ to $\delta'$ using
\begin{equation}\label{eq:fp}
\delta' = 1+f_p^*(-\e^{\ep'}),
\end{equation}
which is proved in \Cref{app:property}. The remainder of the proof is devoted to showing that $(\epsilon', \delta')$-DP guarantees for all $\ep' \ge 0$ is equivalent to the $C_p(f)$-DP guarantee.

%the conversion of this collection back to an $\tilde{f}$-DP guarantee.
%We provide some intuition why the convex combination $f_p = pf+(1-p)\Id$ arises in the main body after the proof.
%, we cannot naively use the ``if and only if'' claim and convert it back to $f_p$-DP, because

At first glance, \eqref{eq:fp} seems to enable the use of \Cref{prop:ftoDP}. Unfortunately, that would be invalid because $f_p$ is asymmetric. To this end, we need to extend \Cref{prop:ftoDP} to general trade-off functions. To avoid conflicting notation, let $g$ be a generic trade-off function, not necessarily symmetric. Denote by $\bar{x}$ be the smallest point such that $g'(x)=-1$, that is, $\bar{x} = \inf\{x\in[0,1]:g'(x)=-1\}$.\footnote{For simplicity, the proof assumes differentiable trade-off functions. If $g$ is not differentiable, use the definition $\bar{x} = \inf\{x\in[0,1]: -1 \in \partial g(x) \}$ instead. This adjustment applies to other parts of the proof.} As a special instance of \Cref{prop:asymm_env} in the appendix, the following result serves our purpose.
\begin{proposition}\label{prop:asymm_for_proof}
If $g(\bar{x})\geqslant\bar{x}$ and a mechanism $M$ is $(\ep,1+g^*(-\e^{\ep}))$-DP for all $\ep\geqslant 0$, then $M$ is $\min\{g,g^{-1}\}^{**}$-DP.
%%\mathrm{Conv}
\end{proposition}


The proof of the present theorem would be complete if \Cref{prop:asymm_for_proof} can be applied to the collection of privacy guarantees in \eqref{eq:fp}for $f_p$. To use \Cref{prop:asymm_for_proof}, it suffices to verify the condition $f_p(\bar{x})\geqslant \bar{x}$ where $\bar{x}$ is the smallest point such that $f_p'(x)=-1$. Let $x^*$ be the (unique) fixed point of $f$. To this end, we collect a few simple facts:
	\begin{itemize}
		% \setlength\itemsep{0.1em}
		\item First, $f'(x^*) = -1$. This is because the graph of $f$ is symmetric with respect to the $45^\circ$ line passing through the origin.
%This is because $f$ is symmetric, so $f = f^{-1}$ have the same derivative at $x^*$. On the other hand, the two derivatives at $x^*$ must also be reciprocal, so it is $\pm 1$. In addition, $f$ is decreasing, so it has to be $-1$.
		\item Second, $\bar{x}\leqslant x^*$. This is because $f_p'(x^*)=pf'(x^*)+(1-p)\Id'(x^*)=-1$ and, by definition, $\bar{x}$ can only be smaller.
%		\item $f_p(x^*)\geqslant f(x^*)$. This is because $f$ is a trade-off function so by \Cref{prop:trade-off} $f\leqslant \Id$, and hence their average $f_p$ lies in between.
	\end{itemize}
With these facts in place, we get
\[f_p(\bar{x})\geqslant f_p(x^*) \geqslant f(x^*) = x^* \geqslant \bar{x}\]
by recognizing that $f_p$ is decreasing and $f_p \ge f$. Hence, the proof is complete.

% So the condition is satisfied and we can apply \Cref{prop:asymm_for_proof} with $g=f_p$, which converts the collection \eqref{eq:fp} into $\mathrm{Conv}(\min\{f_p,f_p^{-1}\})$-DP, i.e. $C_p(f)$-DP.
\end{proof}

% \paragraph{Why $f_p$ Arises in \eqref{eq:fp}} % (fold)
% \label{par:why1}

% In previous literature, $\ep' = \log(1-p + pe^\ep)$ is often approximated by $\log(1-p + pe^\ep)\leqslant -p + pe^\ep =O(p\ep)$. Though convenient in terms of $\ep$, the approximation hides away a nice geometric interpretation:
% $$f_{\ep',\delta'}(x) = pf_{\ep,\delta}(x)+(1-p)\Id(x) \text{ for small } x.$$
% To see this, notice that for small $x$, all three functions involved are linear. Let $k(\cdot)$ and $b(\cdot)$ be the corresponding slope and intercept of the functions. We have
% 	$k(f_{\ep,\delta}) = -e^{\ep} , b(f_{\ep,\delta}) = 1-\delta$ and
% 	$k(\Id)=-1$  $b(\Id)=1$.
% \begin{align*}
% 	k(f_{\ep',\delta'}) &= -e^{\ep'} = -(1-p)-pe^\ep = (1-p)k(\Id)+pk(f_{\ep,\delta})\\
% 	b(f_{\ep',\delta'}) &= 1-\delta' = 1-p\delta = (1-p)+p(1-\delta) = (1-p)b(\Id)+pb(f_{\ep,\delta})
% \end{align*}
% Hence both slope and intercept are averages. The linear function itself must also be an average.

% $\ep$ and $\delta$ are related by $\delta(\ep) = 1+f^*(-e^{\ep})$ so that $f_{\ep,\delta}$ is tangent to $f$. The average $f_{\ep',\delta'} = pf_{\ep,\delta}+(1-p)\Id(x)$ must also be tangent to the average $f_p = pf+(1-p)\Id$. Hence $\ep'$ and $\delta'$ are related by a similar expression $\delta'(\ep') = 1+f_p^*(-e^{\ep})$. This explains why $f_p$ arises in \eqref{eq:fp}.

% \paragraph{Why Symmetrization Induces $\mathrm{Conv}(\min\{g,g^{-1}\})$} % (fold)
% \label{par:why_symmetrization_induces}
% Just like \Cref{prop:DPtof}, we need to find the enveloping function of the collection $f_{\ep,\delta}$ with $\delta=1+g^*(-e^{\ep}),\ep\geqslant0$. Each enveloping $f_{\ep,\delta}$ is in fact the max of two linear functions that are symmetric to each other, one with slope $-e^{\ep}$ and the other $-e^{-\ep}$. We refer to these linear functions as the steeps and flats. The envelope of $f_{\ep,\delta}$ can be obtained by taking the max of the envelopes of steeps and flats. Each $f_{\ep,\delta}$ recovers a point on $g$ with derivative $-e^\ep\leqslant-1$. Hence the collection of steeps envelopes $g_\mathrm{env}$ in the right panel of \Cref{fig:symm_main}. Since flats are symmetric to steeps, they envelope the symmetric function, i.e. $g_\mathrm{env}^{-1}$. From \Cref{fig:symm_main}, it is easy to see that $\max\{g_\mathrm{env},g_\mathrm{env}^{-1}\}$ agrees with $\mathrm{Conv}(\min\{g,g^{-1}\})$.
% paragraph why_symmetrization_induces (end)


% \begin{figure}[t]
%   \centering
%   \includegraphics[width=0.7\linewidth]
% {./figures/Symm_main.pdf}
%   \captionof{figure}{\djs{TODO}}
%   \label{fig:symm_main}
% \end{figure}


% subsection proof_of_subsample_theorems (end)

% paragraph why_ (end)
% Recall in \Cref{thm:subsample} we conclude that the subsample mechanism $\tm$ is $C_p(f)$-DP with
% \[
% C_p(f)(x)=
% 	\left\{
% 	\begin{array}{ll}
% 	f_p(x),&x\in[0,x^*], \\
% 	2px^*+1-p-x, & x\in[x^*,px^*+(1-p)(1-x^*)],\\
% 	f_p^{-1}(x), &x\in[px^*+(1-p)(1-x^*),1]
% 	\end{array}
% 	\right.
% \]
% and $f_p = pf+(1-p)\Id$.
% The proof strategy is to use primal-dual conversion to show that $\tm$ is $\Symm(f_p)$-DP, and then verify that $C_p(f) = \Symm(f_p)$.

%%% Local Variables:
%%% mode: latex
%%% TeX-master: "../paper"
%%% End:

\input{SGD}
\input{discuss}

\subsection*{Acknowledgements}
\label{sec:acknowledgements}

This work was supported in part by NSF grants CCF-1763314 and CNS-1253345, and by the Sloan Foundation, and by a sub-contract for the DARPA Brandeis program.

% This work was supported in part by the NSF via grant CCF-1763314.

{\small
\bibliographystyle{alpha}
\bibliography{privacy}
}


\clearpage
\appendix
\addcontentsline{toc}{section}{Appendices}
% \renewcommand{\thesubsection}{\Alph{subsection}}

%commenting out since there is a compile bug
\input{appendices/appendix_fDP}
\input{appendices/divergence}
\input{appendices/self_contained}
\input{appendices/CLT}
\input{appendices/fastproof}
\input{appendices/property}
\input{appendices/application}







\end{document}

%%% Local Variables:
%%% mode: latex
%%% TeX-master: t
%%% End:
